\section{A hard balanced partitioning core}\label{sec:hardness}

This section isolates a constrained form of structural entropy that is already
computationally hard. The result also explains why cut quality alone does not
characterize a general two-dimensional optimum: module volumes enter the
objective unless they are controlled.

We use the following decision form. An instance consists of a simple cubic
graph $G=(V,E)$ with even $n=|V|$ and an integer $b$. The question is whether
there exists $S\subset V$ with $|S|=n/2$ such that the height-two tree induced
by $\{S,V\setminus S\}$ has entropy at most
$\log_2(n/2)+b/|E|$. The threshold is represented by the integer pair
$(n,b)$; no numerical approximation to the logarithm is part of the input.

For a partition $\mathcal P=\{S_1,\ldots,S_K\}$, let $e_i$ be the internal
edge weight of $S_i$. Directly expanding Eq.~\eqref{eq:tree-se-journal} gives
\begin{align}
 H^{\mathcal P}(G)
 &= -\frac{1}{M}\sum_{u\in V}d_u\log_2\frac{d_u}{V_{S(u)}}
    -\frac{1}{M}\sum_{i=1}^{K}g_i\log_2\frac{V_i}{M}\notag\\
 &= C_G+\frac{1}{M}\sum_{i=1}^{K}
 \left[2e_i\log_2 V_i+g_i\log_2 M\right],
 \label{eq:partition-form}
\end{align}
where $C_G=-M^{-1}\sum_ud_u\log_2d_u$ does not depend on the partition and
$V_i-g_i=2e_i$.

\begin{lemma}[Balanced regular identity]\label{lem:bisection-identity}
Let $G$ be a loop-free $d$-regular graph with even order $n$ and
$m=|E|=dn/2$. For a balanced bipartition $\mathcal P$ with cut size
$c(\mathcal P)$,
\begin{equation}
 H^{\mathcal P}(G)=\log_2\frac{n}{2}+\frac{c(\mathcal P)}{m}.
 \label{eq:balanced-identity}
\end{equation}
Consequently, for two balanced bipartitions $\mathcal P$ and $\mathcal Q$,
$H^{\mathcal P}(G)-H^{\mathcal Q}(G)=
[c(\mathcal P)-c(\mathcal Q)]/m$.
\end{lemma}

\begin{proof}
Every vertex has degree $d$, both modules have volume $dn/2=M/2$, and both
module cuts equal $c(\mathcal P)$. The leaf contribution to the height-two
tree is
\[
 -\frac{1}{M}\sum_{u\in V}d\log_2\frac{d}{M/2}
 =\log_2(n/2).
\]
The two module contributions sum to
$-2c(\mathcal P)M^{-1}\log_2(1/2)=c(\mathcal P)/m$.
Adding the terms proves Eq.~\eqref{eq:balanced-identity}.
\end{proof}

\begin{theorem}[Constrained structural-entropy hardness]
\label{thm:balanced-hardness}
The decision form of balanced two-module structural entropy is NP-complete on
simple unweighted cubic graphs. Consequently, computing an optimal balanced
two-module encoding tree is NP-hard on the same class.
\end{theorem}

\begin{proof}
Minimum Bisection is NP-complete on simple cubic graphs
\cite{bui1987graph}. Map an instance $(G,b)$ identically to the entropy
decision instance $(G,b)$. The mapping is polynomial and introduces neither
weights nor auxiliary vertices. By Lemma~\ref{lem:bisection-identity}, a
balanced partition $\mathcal P$ satisfies
$H^{\mathcal P}(G)\leq\log_2(n/2)+b/m$ if and only if
$c(\mathcal P)\leq b$. This proves NP-hardness. A balanced set $S$ is a
polynomial certificate, and the equivalent integer comparison
$c(\mathcal P)\leq b$ verifies it in polynomial time; hence the decision
problem is in NP. An optimizer would decide the problem by comparing its
returned cut size with $b$, so the optimization form is NP-hard.
\end{proof}

The balance constraint cannot simply be discarded. On the 10-vertex
pentagonal prism and M\"obius ladder, exhaustive enumeration finds a $4/6$
partition with entropy $2.625497$ bits, below the best balanced value
$2.655261$. Thus regularity and exactly two nonempty modules do not force a
bisection. This counterexample delineates the theorem's scope and rules out an
otherwise tempting transfer to unconstrained two-module optimization.
